% =====================================================================
%  Lecture 15 -- Software-Defined Networking (SDN)         (Persian edition)
% =====================================================================
\documentclass[11pt,a4paper]{article}
\providecommand\ACNroot{../..}
\input{\ACNroot/templates/lecture-fa}

\ACNset{number}{15}
\ACNset{title}{شبکه‌سازی نرم‌افزارمحور (\lr{SDN})}
\ACNset{shorttitle}{شبکه‌سازی نرم‌افزارمحور}
\ACNset{titleother}{Software-Defined Networking (SDN)}
\ACNset{part}{بخش سوم -- معماری‌های نرم‌افزارمحور (درس‌های ۱۵ تا ۱۸)}
\ACNset{time}{۶ تا ۸ ساعت (۳ ساعت مطالعه، ۳ تا ۵ ساعت تمرین)}
\ACNset{level}{متوسط}
\ACNset{prereq}{مسیریابی و سوئیچینگ \lr{IP}؛ درس ۱۳ (\lr{CUPS})؛ صف \lr{M/M/1} (درس ۱۲)}
\ACNset{gen}{5}

\begin{document}
\ACNcover

% ---------------------------------------------------------------------
\section{اهداف یادگیری}
پس از مطالعهٔ این درس باید بتوانید:
\begin{objectives}
  \item جداسازی صفحه‌های کنترل و داده و انگیزهٔ کنترل منطقاً متمرکز را \textbf{توضیح دهید}.
  \item لایه‌ها و واسط‌های \acro{SDN} (شمالی، جنوبی) و مدل سوئیچ \lr{OpenFlow} -- جدول‌های جریان، خط لوله، گروه‌ها و
        سنجه‌ها -- را \textbf{توصیف کنید}.
  \item نصب جریان واکنشی و پیش‌کنشی را \textbf{دنبال کنید} و تأخیر و نیاز جدولی آن‌ها را \textbf{محاسبه کنید}.
  \item برنامه‌پذیری صفحهٔ داده با \lr{P4} را \textbf{توضیح دهید}.
  \item جایابی کنترل‌گر و بار کنترل‌گر را \textbf{تحلیل کنید}.
  \item اصول \acro{SDN} را در شبکه‌های موبایل (\acro{CUPS}، \acro{PFCP}، انتقال و برش‌بندی) \textbf{شناسایی کنید}.
\end{objectives}

% ---------------------------------------------------------------------
\section{دانش پیش‌نیاز}
باید بدانید مسیریاب‌های \acro{IP} و سوئیچ‌های اترنت چگونه بسته‌ها را هدایت می‌کنند (تطابق طولانی‌ترین پیشوند، یادگیری
\acro{MAC}) و اینکه پروتکل‌های مسیریابی مانند \acro{OSPF} و \acro{BGP} جدول‌های هدایت را به‌صورت توزیع‌شده محاسبه می‌کنند. از
این درس‌نامه، به \acro{CUPS} و \acro{PFCP} (درس ۱۳) و فرمول تأخیر \lr{M/M/1} (درس ۱۲) نیاز دارید.

\paragraph{واژه‌های کلیدی} واژگان این درس همراه با معادل انگلیسی آن‌ها:
\begin{keyterms}
صفحهٔ کنترل & \lr{Control plane} & منطقی که تعیین می‌کند ترافیک چگونه هدایت شود. \\
صفحهٔ داده & \lr{Data plane} & هدایت هر بسته بر اساس حالت نصب‌شده. \\
کنترل‌گر \acro{SDN} & \lr{SDN controller} & نرم‌افزار منطقاً متمرکزی که صفحهٔ داده را برنامه‌ریزی می‌کند. \\
\acro{API} جنوبی & \lr{Southbound API} & واسط کنترل‌گر--دستگاه (\lr{OpenFlow}، \lr{P4Runtime}، \lr{NETCONF}). \\
\acro{API} شمالی & \lr{Northbound API} & واسط کنترل‌گر--کاربرد (مثلاً \acro{REST}). \\
جدول جریان & \lr{Flow table} & قواعد تطابق--عمل اولویت‌دار در یک سوئیچ. \\
پیام‌های \lr{OpenFlow} & \lr{Packet-In / Flow-Mod} & گزارش بسته از سوئیچ به کنترل‌گر / نصب قاعده. \\
واکنشی / پیش‌کنشی & \lr{Reactive / proactive} & نصب قواعد با نخستین بسته / از پیش. \\
زبان \lr{P4} & \lr{P4} & زبانی برای برنامه‌ریزی پردازنده‌های بستهٔ مستقل از پروتکل. \\
حافظهٔ آدرس‌پذیر محتوایی سه‌حالته & \lr{TCAM} & حافظهٔ سخت‌افزاری برای تطابق با نویسهٔ عام. \\
\end{keyterms}

% ---------------------------------------------------------------------
\section{مقدمه و انگیزه}
در یک شبکهٔ سنتی هر مسیریاب نرم‌افزار کنترلی خودش را اجرا می‌کند: با همسایگان پیام‌های مسیریابی مبادله می‌کند، مسیرها را
محاسبه و جدول هدایت خود را پر می‌کند. کنترل و هدایت در هر جعبه به هم بسته‌اند و هر فروشنده پروتکل‌ها را در نرم‌افزار
بستهٔ خود پیاده می‌کند. این طراحی مقاوم است، اما تغییر شبکه را دشوار می‌کند: یک رفتار هدایت جدید به پروتکلی جدید،
استانداردی و سال‌ها پیاده‌سازی نیاز دارد، و سیاست‌های سراسری شبکه (مهندسی ترافیک، جداسازی، کنترل دسترسی) باید به‌طور
غیرمستقیم از طریق پیکربندی‌های بسیاری در هر جعبه بیان شوند.

شبکه‌سازی نرم‌افزارمحور (\acro{SDN}) این دو صفحه را جدا می‌کند \cite{mckeown2008openflow,kreutz2015sdn}. دستگاه‌های هدایت به
موتورهای ساده و برنامه‌پذیر تطابق--عمل تبدیل می‌شوند؛ یک کنترل‌گر منطقاً متمرکز با دید سراسری حالت آن‌ها را محاسبه و از طریق
واسطی باز نصب می‌کند. رفتار شبکه به نرم‌افزاری تبدیل می‌شود که می‌توان آن را مانند هر نرم‌افزار دیگری نوشت، آزمود و مستقر
کرد. شبکه‌های موبایل همین ایده را پذیرفتند: \acro{CUPS} و \acro{PFCP} (درس ۱۳) همان \acro{SDN} اعمال‌شده بر دروازه‌های
بسته‌ای‌اند و هستهٔ نسل پنجم با \acro{UPF} آن بر این پایه ساخته شده است. این درس بخش سوم را می‌گشاید که فناوری‌های
نرم‌افزاری -- \acro{SDN}، \acro{NFV} و معماری سرویس‌محور -- را که شبکه‌های نسل پنجم بر آن‌ها استوارند پوشش می‌دهد.

% ---------------------------------------------------------------------
\section{مفاهیم بنیادی}

\subsection{سه اصل}
\begin{enumerate}
  \item \textbf{جداسازی صفحه‌های کنترل و داده}: دستگاه‌ها هدایت می‌کنند؛ کنترل‌گرها تصمیم می‌گیرند.
  \item \textbf{کنترل منطقاً متمرکز}: یک کنترل‌گر (احتمالاً توزیع‌شده و تکثیرشده) دید سراسری از توپولوژی و حالت دارد.
  \item \textbf{برنامه‌پذیری}: واسط‌های باز به کاربردهای نرم‌افزاری اجازه می‌دهند رفتار شبکه را تعریف کنند.
\end{enumerate}

\subsection{لایه‌ها و واسط‌ها}
\lr{RFC 7426} \acro{SDN} را بر حسب صفحه‌ها -- هدایت، عملیاتی، کنترل، مدیریت و کاربرد -- و لایه‌های انتزاع میان آن‌ها توصیف
می‌کند \cite{rfc7426}. در نمای سادهٔ رایج (\cref{fig:arch}) کاربردها از \acro{API} \emph{شمالی} کنترل‌گر (اغلب \acro{REST})
استفاده می‌کنند و کنترل‌گر دستگاه‌ها را از طریق \acro{API} \emph{جنوبی} برنامه‌ریزی می‌کند: \lr{OpenFlow}
\cite{onf2015openflow} برای جدول‌های جریان، \lr{P4Runtime} برای دستگاه‌های برنامه‌ریزی‌شده با \lr{P4} یا \lr{NETCONF/YANG}
برای پیکربندی.

\begin{acnfigure}
  \resizebox{\linewidth}{!}{\acnfig{lecture-15/sdn-architecture}}
  \caption{لایه‌های \acro{SDN}: کاربردها، کنترل‌گر و صفحهٔ داده با \acro{API}های شمالی و جنوبی.}
  \label{fig:arch}
\end{acnfigure}

\subsection{سوئیچ \lr{OpenFlow}}
\begin{definitionbox}{مدخل جریان}
یک مدخل جریان شامل فیلدهای تطابق (درگاه ورودی و فیلدهای سرآیند بسته از اترنت تا لایهٔ انتقال)، یک اولویت، شمارنده‌ها،
دستورها (اعمال یا نوشتن عمل‌ها، رفتن به جدولی دیگر)، مهلت‌های بیکاری و سخت و یک \lr{cookie} مبهم است
\cite{onf2015openflow}. بسته با مدخل‌های یک جدول تطبیق داده می‌شود؛ تطابق با بالاترین اولویت برنده است.
\end{definitionbox}
هر سوئیچ یک خط لوله از جدول‌های جریان (\cref{fig:table})، یک \emph{جدول گروه} (برای چندپخشی، توازن بار با گروه‌های
\emph{select} و جایگزینی سریع) و یک \emph{جدول سنجه} (محدودسازی نرخ) دارد. اگر هیچ مدخلی تطابق نیابد (\emph{عدم تطابق
جدول})، مدخل عدم تطابق تصمیم می‌گیرد: دورریختن، ادامه یا ارسال بسته به کنترل‌گر. سوئیچ از طریق یک اتصال \acro{TCP} (درگاه
۶۶۵۳، با \acro{TLS} اختیاری) با کنترل‌گر گفت‌وگو می‌کند.

\begin{acnfigure}
  \resizebox{\linewidth}{!}{\acnfig{lecture-15/flow-table}}
  \caption{خط لولهٔ \lr{OpenFlow} با چند جدول جریان و فیلدهای یک مدخل جریان.}
  \label{fig:table}
\end{acnfigure}

% ---------------------------------------------------------------------
\section{توضیحات فنی تفصیلی}

\subsection{پیام‌های \lr{OpenFlow}}
\begin{itemize}
  \item \textbf{کنترل‌گر به سوئیچ}: \lr{Features} (قابلیت‌ها)، \lr{Flow-Mod} (افزودن، تغییر، حذف مدخل‌ها)، \lr{Group-Mod}،
        \lr{Meter-Mod}، \lr{Packet-Out} (ارسال یک بسته از یک درگاه)، \lr{Multipart/Stats} (خواندن شمارنده‌ها)، \lr{Barrier}
        (اطمینان از پردازش پیام‌های پیشین).
  \item \textbf{ناهمگام}: \lr{Packet-In} (بسته‌ای که به کنترل‌گر فرستاده می‌شود، به‌طور پیش‌فرض با ۱۲۸ بایت نخست آن و یک
        شناسهٔ بافر)، \lr{Flow-Removed} (یک مدخل منقضی شد)، \lr{Port-Status} (یک درگاه تغییر کرد).
  \item \textbf{متقارن}: \lr{Hello}، \lr{Echo} (زنده‌بودن)، \lr{Error}.
\end{itemize}

\subsection{کنترل واکنشی و پیش‌کنشی}
در حالت \emph{واکنشی} نخستین بستهٔ یک جریان جدید تطابق نمی‌یابد، سوئیچ یک \lr{Packet-In} می‌فرستد و کنترل‌گر مسیری محاسبه و
مدخل‌هایی در طول آن نصب می‌کند و سپس بسته را با \lr{Packet-Out} آزاد می‌کند (\cref{fig:reactive}). کنترل ریزدانه ممکن است،
اما هر جریان جدید هزینهٔ یک رفت‌وبرگشت تا کنترل‌گر را می‌پردازد و بار کنترل‌گر با نرخ ورود جریان رشد می‌کند. در حالت
\emph{پیش‌کنشی} کنترل‌گر قواعد (معمولاً با نویسهٔ عام) را از پیش و بر اساس دانش خود از توپولوژی و سیاست نصب می‌کند؛ هیچ بسته‌ای
منتظر نمی‌ماند و کنترل‌گر خارج از مسیر بحرانی است، اما قواعد باید تجمیع شوند تا در جدول‌ها جا شوند. استقرارهای عملی هر دو را
ترکیب می‌کنند: قواعد پیش‌کنشی برای بیشتر ترافیک و رسیدگی واکنشی برای استثناها.

\begin{acnfigure}
  \resizebox{\linewidth}{!}{\acnfig{lecture-15/reactive-flow}}
  \caption{برقراری واکنشی جریان: عدم تطابق جدول، \lr{Packet-In}، \lr{Flow-Mod} در طول مسیر، \lr{Packet-Out}.}
  \label{fig:reactive}
\end{acnfigure}

\subsection{برنامه‌پذیری صفحهٔ داده: \lr{P4}}
\lr{OpenFlow} فیلدهای سرآیندی را که سوئیچ می‌تواند تطبیق دهد ثابت می‌کند؛ افزودن یک پروتکل جدید به نسخهٔ جدیدی از مشخصات و
سخت‌افزار نیاز دارد. \lr{P4} \cite{bosshart2014p4} یک گام فراتر می‌رود: برنامه سرآیندها، ماشین حالت تجزیه‌گر، جدول‌های
تطابق--عمل و جریان کنترل را تعریف می‌کند و یک کامپایلر آن را بر یک هدف برنامه‌پذیر (تراشهٔ سوئیچ، کارت شبکهٔ هوشمند، سوئیچ
نرم‌افزاری) نگاشت می‌کند که به‌صورت یک تجزیه‌گر، خطوط لولهٔ تطابق--عمل ورودی و خروجی و یک وارون‌تجزیه‌گر سازمان یافته است
(\cref{fig:p4}). سپس کنترل‌گر جدول‌ها را در زمان اجرا از طریق \lr{P4Runtime} پر می‌کند. برای نمونه، از \lr{P4} برای
پیاده‌سازی کپسوله‌سازی \lr{GTP-U} و کارکردهای \acro{UPF} مستقیماً در سخت‌افزار سوئیچ استفاده شده است.

\begin{acnfigure}
  \resizebox{\linewidth}{!}{\acnfig{lecture-15/p4-pipeline}}
  \caption{معماری سوئیچ مستقل از پروتکل که با \lr{P4} برنامه‌ریزی می‌شود.}
  \label{fig:p4}
\end{acnfigure}

\subsection{کنترل‌گرهای توزیع‌شده و سازگاری}
یک کنترل‌گر واحد نقطهٔ شکست و گلوگاه است. کنترل‌گرهای عملیاتی به‌صورت خوشه‌هایی اجرا می‌شوند که حالت را با پروتکل‌های
اجماع تکثیر و سوئیچ‌ها را میان نمونه‌ها تقسیم می‌کنند؛ هر سوئیچ با نقش‌های اصلی/فرعی به چند کنترل‌گر وصل می‌شود.
به‌روزرسانی بسیاری از سوئیچ‌ها اتمی نیست: در طول یک به‌روزرسانی برخی بسته‌ها ممکن است قواعد قدیم را در یک سوئیچ و قواعد
جدید را در سوئیچ بعدی ببینند و حلقه یا نقض سیاست پدید آید. روش‌های \emph{به‌روزرسانی سازگار} بسته‌ها را در ورودی با نسخهٔ
پیکربندی برچسب می‌زنند و پیش از تغییر برچسب‌ها نسخهٔ جدید را روی همهٔ سوئیچ‌ها نصب می‌کنند، به‌طوری که هر بسته یا
پیکربندی قدیم یا جدید را ببیند و هرگز آمیزه‌ای از آن‌ها را نبیند.

\subsection{\lr{SDN} در شبکه‌های موبایل}
\begin{itemize}
  \item \textbf{دروازه‌ها}: \acro{CUPS} (درس ۱۳) کنترل \lr{SGW/PGW} را از توابع صفحهٔ کاربر جدا می‌کند؛ \acro{PFCP} نقش یک
        پروتکل جنوبی با قواعدی (\lr{PDR, FAR, QER, URR}) شبیه مدخل‌های \lr{OpenFlow} را بازی می‌کند.
  \item \textbf{هستهٔ نسل پنجم}: \acro{SMF} از طریق \lr{N4} با \acro{PFCP} \acro{UPF}ها را کنترل می‌کند (درس ۲۳).
  \item \textbf{انتقال}: کنترل‌گرهای \acro{SDN} مسیرهای بک‌هال و فرانت‌هال، مسیریابی قطعه‌ای و برش‌های انتقال را پیکربندی
        می‌کنند.
  \item \textbf{\acro{RAN}}: معماری \lr{O-RAN} کنترل \acro{RAN} را از طریق یک کنترل‌گر هوشمند \acro{RAN} به روی کاربردهای
        نرم‌افزاری می‌گشاید -- ایدهٔ \acro{SDN} اعمال‌شده بر مدیریت منابع رادیویی.
\end{itemize}

% ---------------------------------------------------------------------
\section{مدل‌های ریاضی}

\subsection{تأخیر برقراری واکنشی جریان}\label{sec:setup}
\paragraph{مدل} نخستین بستهٔ یک جریان منتظر رفت‌وبرگشت سوئیچ--کنترل‌گر $\tau$ و پردازش کنترل‌گر می‌ماند. با $S$ سوئیچ که هر
کدام $\lambda_s$ جریان جدید در ثانیه تولید می‌کنند و کنترل‌گری که به‌صورت یک خدمت‌دهندهٔ \lr{M/M/1} با نرخ $\mu$ مدل شده،
\begin{equation}\label{eq:setup}
  T_{\mathrm{setup}} \approx \tau + \frac{1}{\mu - S\lambda_s}, \qquad S\lambda_s < \mu .
\end{equation}
\paragraph{مثال} $\tau = 2$ میلی‌ثانیه، $S = 100$، $\lambda_s = 1000$\,s$^{-1}$ و $\mu = 1.5\times10^5$\,s$^{-1}$:
$T_{\mathrm{setup}} = 2 + 0.02 = 2.02$ میلی‌ثانیه. در $S\lambda_s = 1.49\times10^5$ جملهٔ صف به ۱ میلی‌ثانیه می‌رسد
(\cref{fig:setup}).
\paragraph{تفسیر} در بار متوسط تأخیر انتشار تا کنترل‌گر غالب است، که استدلالی برای کنترل‌گرهای نزدیک به سوئیچ‌هاست؛ نزدیک
اشباع جملهٔ صف منفجر می‌شود و طوفان‌های \lr{Packet-In} (مثلاً هنگام یک پویش یا خرابی) می‌توانند کنترل‌گر را از کار
بیندازند -- یک دلیل برای ترجیح قواعد پیش‌کنشی.

\begin{acnfigure}
  \acnfig{lecture-15/setup-delay}
  \caption{تأخیر صف کنترل‌گر بر حسب نرخ \lr{Packet-In} (\cref{eq:setup}).}
  \label{fig:setup}
\end{acnfigure}

\subsection{اشغال جدول جریان}\label{sec:little}
\paragraph{مدل} با قواعد تطابق دقیق (برای هر جریان)، جریان‌های جدید با نرخ $\lambda$ می‌رسند و هر مدخل به اندازهٔ طول جریان
به‌علاوهٔ مهلت بیکاری، به‌طور میانگین $\bar T$، می‌ماند. بنا به قانون لیتل تعداد میانگین مدخل‌ها
\begin{equation}\label{eq:little}
  \bar N = \lambda\,\bar T
\end{equation}
است.
\paragraph{مثال} $\lambda = 2000$\,s$^{-1}$، طول میانگین جریان ۱۰ ثانیه و مهلت بیکاری ۱۰ ثانیه: $\bar N = 40\,000$ مدخل --
بسیار بیشتر از چند هزار مدخل \acro{TCAM} بسیاری از سوئیچ‌ها. کنترل واکنشی برای هر جریان در چنین سوئیچ‌هایی مقیاس‌پذیر
نیست؛ قواعد با نویسهٔ عام، تجمیع یا جدول‌های تطابق دقیق بزرگ‌تر (مبتنی بر درهم‌سازی) لازم‌اند.

\subsection{جایابی کنترل‌گر}\label{sec:place}
\paragraph{مدل} برای سوئیچ‌های $v\in V$ و مکان‌های کنترل‌گر $C$، با $d(v,c)$ به‌عنوان تأخیر میان آن‌ها، تأخیر میانگین و بدترین
حالت سوئیچ--کنترل‌گر عبارت‌اند از
\begin{equation}\label{eq:place}
  \bar L(C) = \frac{1}{|V|}\sum_{v\in V}\min_{c\in C} d(v,c), \qquad
  L_{\max}(C) = \max_{v\in V}\min_{c\in C} d(v,c).
\end{equation}
کمینه‌سازی $\bar L$ برای $|C| = k$ مسئلهٔ $k$-میانه و کمینه‌سازی $L_{\max}$ مسئلهٔ $k$-مرکز است.
\paragraph{مثال} پنج سوئیچ روی یک خط با فاصلهٔ ۱ میلی‌ثانیه: یک کنترل‌گر در مرکز $\bar L = 1.2$ میلی‌ثانیه ($L_{\max} = 2$
میلی‌ثانیه) می‌دهد؛ در یک انتها $\bar L = 2$ میلی‌ثانیه. دو کنترل‌گر در سوئیچ‌های دوم و چهارم $\bar L = 0.6$ میلی‌ثانیه و
$L_{\max} = 1$ میلی‌ثانیه می‌دهند.

% ---------------------------------------------------------------------
\section{معماری شبکه و نمودارها}
در یک مرکز داده، کنترل‌گر \acro{SDN} تار برگ--ستون و سوئیچ‌های مجازی درون کارسازها را برنامه‌ریزی می‌کند؛ در یک شبکهٔ گسترده،
مسیرهای مهندسی‌شدهٔ ترافیک میان سایت‌ها را محاسبه می‌کند؛ در یک اپراتور موبایل، کنترل‌گرهای \acro{SDN} شبکه‌های انتقال را
مدیریت می‌کنند و از طریق توابع کنترلی هسته، گره‌های صفحهٔ کاربر را. خود کنترل‌گر یک کاربرد توزیع‌شده است که در نمونه‌های
افزونه و اغلب روی زیرساخت مجازی‌شدهٔ درس ۱۶ مستقر می‌شود.

% ---------------------------------------------------------------------
\section{رویه‌های گام‌به‌گام}

\begin{procedurebox}{نصب واکنشی جریان}\label{proc:reactive}
\begin{steps}
  \item نخستین بستهٔ یک جریان جدید به سوئیچ ۱ می‌رسد و با هیچ مدخلی تطابق نمی‌یابد؛ مدخل عدم تطابق آن را به‌صورت
        \lr{Packet-In} (سرآیند و شناسهٔ بافر) به کنترل‌گر می‌فرستد.
  \item کنترل‌گر مبدأ و مقصد را شناسایی، سیاست را بررسی و با دید توپولوژی خود مسیری محاسبه می‌کند.
  \item پیام‌های \lr{Flow-Mod} را به همهٔ سوئیچ‌های مسیر می‌فرستد (نخست سمت مقصد، تا از عدم تطابق‌های بیشتر پرهیز شود) و
        به‌طور اختیاری یک \lr{Barrier} پس از آن.
  \item \lr{Packet-Out} را به سوئیچ ۱ می‌فرستد تا بستهٔ بافرشده در مسیر جدید آزاد شود.
  \item بسته‌های بعدی با مدخل‌های جدید تطابق می‌یابند و با نرخ خط هدایت می‌شوند؛ مدخل‌ها پس از مهلت بیکاری منقضی می‌شوند و
        \lr{Flow-Removed} شمارنده‌های آن‌ها را گزارش می‌دهد.
\end{steps}
\end{procedurebox}

\begin{procedurebox}{کشف توپولوژی}\label{proc:lldp}
\begin{steps}
  \item هنگام اتصال، هر سوئیچ درگاه‌های خود را گزارش می‌دهد (پاسخ \lr{Features}).
  \item کنترل‌گر از طریق \lr{Packet-Out} یک قاب کشف پیوند (\acro{LLDP}) روی هر درگاه هر سوئیچ می‌فرستد.
  \item سوئیچ همسایه قاب را دریافت می‌کند و با یک قاعدهٔ ازپیش‌نصب‌شده آن را همراه با درگاه ورودی‌اش به‌صورت \lr{Packet-In}
        برمی‌گرداند.
  \item کنترل‌گر از (سوئیچ و درگاه فرستنده) و (سوئیچ و درگاه گیرنده) هر پیوند را می‌آموزد؛ این کار را دوره‌ای تکرار و توپولوژی
        را با رویدادهای \lr{Port-Status} به‌روز می‌کند.
\end{steps}
\end{procedurebox}

% ---------------------------------------------------------------------
\section{مثال‌های حل‌شده}

\begin{examplebox}{بودجهٔ تأخیر برقراری}\label{ex:setup}
\textbf{مسئله.} یک شبکهٔ پردیس ۱۰۰ سوئیچ دارد؛ هر کدام ۱۰۰۰ جریان جدید در ثانیه می‌بیند. کنترل‌گر $1.5\times10^5$
\lr{Packet-In} در ثانیه را پردازش می‌کند و ۱ میلی‌ثانیه (یک‌طرفه) فاصله دارد. نخستین بستهٔ یک جریان چه تأخیری می‌بیند و این چه
کسری از جریان‌های کوتاه وب (۱۰ بسته، ۲۰ میلی‌ثانیه) است؟

\textbf{حل.} $\tau = 2$ میلی‌ثانیه و صف ۰٫۰۲ میلی‌ثانیه (\cref{eq:setup}): حدود ۲ میلی‌ثانیه، یعنی ۱۰٪ طول جریان. کنترل‌گر با
۶۷٪ بار کار می‌کند؛ افزایش ۵۰ درصدی نرخ جریان آن را اشباع می‌کند.
\end{examplebox}

\begin{examplebox}{اندازهٔ جدول}\label{ex:table}
\textbf{مسئله.} یک سوئیچ بالای قفسه با ۴۰۰۰ مدخل \acro{TCAM}، ۲۰۰۰ جریان جدید در ثانیه می‌بیند که هر کدام به‌طور میانگین ۱۰
ثانیه طول می‌کشد و مهلت بیکاری ۱۰ ثانیه است. آیا می‌توان از قواعد واکنشی برای هر جریان استفاده کرد؟ اگر مهلت بیکاری ۱ ثانیه
باشد چه؟

\textbf{حل.} $\bar N = 2000\times20 = 40\,000 > 4000$: خیر. با مهلت بیکاری ۱ ثانیه، $\bar N = 22\,000$: باز هم خیر. قواعد باید
تجمیع شوند، مثلاً برای هر پیشوند مقصد (پیش‌کنشی)، که به مدخل برای هر مقصد نیاز دارد نه برای هر جریان.
\end{examplebox}

\begin{examplebox}{جایابی}\label{ex:place}
\textbf{مسئله.} مقادیر مثال خطی \cref{sec:place} را بررسی کنید.

\textbf{حل.} مرکز (سوئیچ ۳): فاصله‌های ۲، ۱، ۰، ۱، ۲؛ میانگین ۱٫۲ و بیشینه ۲ میلی‌ثانیه. انتها (سوئیچ ۱): ۰، ۱، ۲، ۳، ۴؛ میانگین ۲ و
بیشینه ۴ میلی‌ثانیه. کنترل‌گرها در ۲ و ۴: ۱، ۰، ۱، ۰، ۱؛ میانگین ۰٫۶ و بیشینه ۱ میلی‌ثانیه.
\end{examplebox}

% ---------------------------------------------------------------------
\section{مقایسهٔ فناوری‌ها}
\begin{table}[htbp]
\centering\small
\caption{کنترل توزیع‌شدهٔ سنتی در برابر \acro{SDN}.}\label{tab:cmp}
\rowcolors{2}{acnmist}{white}
\begin{tabularx}{\linewidth}{@{}l Y Y@{}}
\toprule
\rowcolor{white} & کنترل توزیع‌شده (\acro{OSPF}، \acro{BGP}) & \acro{SDN} \\
\midrule
محل کنترل        & در هر دستگاه                          & کنترل‌گر منطقاً متمرکز \\
دید              & محلی، ساخته‌شده با تبادل پروتکل        & توپولوژی و حالت سراسری \\
رفتار جدید       & پروتکل، استاندارد و میان‌افزار جدید    & کاربرد جدید کنترل‌گر \\
رسیدگی به خرابی  & محلی، سریع، خودمختار                  & به کنترل‌گر یا پشتیبان‌های ازپیش‌نصب‌شده نیاز دارد \\
حد مقیاس‌پذیری   & همگرایی پروتکل، اندازهٔ جدول           & ظرفیت کنترل‌گر، تأخیر کانال کنترل \\
بیان سیاست       & غیرمستقیم (معیارها، فیلترها در هر جعبه) & مستقیم، سراسری \\
وابستگی به فروشنده & سخت‌افزار و نرم‌افزار به هم بسته      & واسط‌های باز، تراشه‌های تجاری \\
\bottomrule
\end{tabularx}
\end{table}

\paragraph{نه جایگزینی برای مسیریابی} \acro{SDN} پروتکل‌های توزیع‌شده را منسوخ نمی‌کند: بسیاری از استقرارها \acro{BGP} یا
\lr{IS-IS} را برای دسترس‌پذیری پایه و جایگزینی سریع محلی نگه می‌دارند و کنترل‌گر را برای سیاست و بهینه‌سازی روی آن به کار
می‌برند (\acro{SDN} ترکیبی).

% ---------------------------------------------------------------------
\section{نکته‌های کلیدی}
\begin{takeawaybox}
\begin{itemize}
  \item \acro{SDN} کنترل را از هدایت جدا، کنترل را منطقاً متمرکز و شبکه را برنامه‌پذیر می‌کند.
  \item سوئیچ‌های \lr{OpenFlow}: جدول‌های جریان اولویت‌دار در خط لوله، گروه‌ها، سنجه‌ها؛ \lr{Packet-In}، \lr{Flow-Mod}،
        \lr{Packet-Out}.
  \item کنترل واکنشی برای هر جریان جدید $T_{\mathrm{setup}} \approx \tau + 1/(\mu - S\lambda_s)$ هزینه دارد؛ کنترل پیش‌کنشی
        به قواعد تجمیع‌شده نیاز دارد.
  \item اشغال جدول از قانون لیتل پیروی می‌کند، $\bar N = \lambda\bar T$؛ قواعد هر جریان به‌ندرت در \acro{TCAM} جا می‌شوند.
  \item \lr{P4} تجزیه‌گر و خط لوله را برنامه‌پذیر می‌کند؛ جایابی کنترل‌گر یک مسئلهٔ $k$-میانه/$k$-مرکز است.
  \item \lr{CUPS/PFCP} و \acro{UPF} نسل پنجم \acro{SDN} را بر هستهٔ موبایل اعمال می‌کنند.
\end{itemize}
\end{takeawaybox}

% ---------------------------------------------------------------------
\section{تمرین‌ها}

\begin{exercisebox}{1}{کدام صفحه؟}\label{exr:planes}
به صفحهٔ کنترل یا داده نسبت دهید: (الف) محاسبهٔ کوتاه‌ترین مسیر \acro{OSPF}؛ (ب) کاهش \acro{TTL}؛ (ج) نصب یک \lr{Flow-Mod}؛
(د) تطبیق یک بسته با \acro{TCAM}؛ (هـ) کشف توپولوژی مبتنی بر \acro{LLDP}.
\end{exercisebox}

\begin{exercisebox}{1}{واکنشی یا پیش‌کنشی؟}\label{exr:proactive}
برای هر مورد کنترل واکنشی یا پیش‌کنشی را برگزینید و توجیه کنید: (الف) تار یک مرکز داده با هدایت بر اساس مقصد؛ (ب) شبکهٔ
پردیسی که میهمانان پیش از مجازشدن نخستین جریان خود باید احراز هویت شوند؛ (ج) شبکهٔ بک‌هال موبایل.
\end{exercisebox}

\begin{exercisebox}{2}{پهنای باند کانال کنترل}\label{exr:packetin}
یک کنترل‌گر $10^5$ پیام \lr{Packet-In} در ثانیه دریافت می‌کند که هر کدام ۱۲۸ بایت از بسته به‌علاوهٔ ۳۲ بایت سربار
\lr{OpenFlow} و \lr{TCP/IP} (سرشکن‌شده) دارند. پهنای باند کانال کنترل را محاسبه کنید. اگر سوئیچ به‌جای بافرکردن، کل
بسته‌های ۱۵۰۰ بایتی را بفرستد چه تغییری می‌کند؟
\end{exercisebox}

\begin{exercisebox}{2}{مهلت‌ها}\label{exr:little}
با \cref{eq:little} مهلت بیکاری‌ای را بیابید که اشغال میانگین را برای ۱۰۰۰ جریان جدید در ثانیه با طول میانگین ۸ ثانیه در
۱۵\,۰۰۰ مدخل نگه دارد. بهای یک مهلت بیکاری بسیار کوتاه چیست؟
\end{exercisebox}

\begin{exercisebox}{3}{جایابی روی حلقه}\label{exr:place}
شش سوئیچ یک حلقه با پیوندهای ۱ میلی‌ثانیه‌ای می‌سازند. $\bar L$ و $L_{\max}$ را برای یک کنترل‌گر و برای دو کنترل‌گر با جایابی
بهینه محاسبه کنید. اگر یکی از دو کنترل‌گر خراب شود بدترین حالت چه می‌شود؟
\end{exercisebox}

\begin{exercisebox}{3}{به‌روزرسانی سازگار}\label{exr:consistent}
یک تغییر مسیر ترافیک را از مسیر \lr{A--B--D} به \lr{A--C--D} می‌برد. نشان دهید به‌روزرسانی سوئیچ \lr{A} پیش از نصب قواعد
روی \lr{C} چگونه می‌تواند بسته‌ها را دور بریزد، و یک به‌روزرسانی با برچسب نسخه را توصیف کنید که هم از دورریختن و هم از آمیختن
پرهیز کند. این روش به‌طور موقت به چند قاعده نیاز دارد؟
\end{exercisebox}

\subsection*{پاسخ‌ها}
\addcontentsline{toc}{subsection}{پاسخ‌ها}

\begin{solutionbox}{\ref{exr:planes}}
(الف) کنترل؛ (ب) داده؛ (ج) کنترل (عمل نصب؛ مدخل سپس توسط صفحهٔ داده استفاده می‌شود)؛ (د) داده؛ (هـ) کنترل.
\end{solutionbox}

\begin{solutionbox}{\ref{exr:proactive}}
(الف) پیش‌کنشی: مقصدها معلوم‌اند، حجم ترافیک عظیم است و تأخیر برقراری برای هر جریان پذیرفتنی نیست. (ب) واکنشی برای نخستین
جریان هر میهمان (کنترل‌گر باید آن را ببیند تا احراز هویت را اعمال کند)، سپس قواعد پیش‌کنشی برای میزبان‌های مجاز. (ج) پیش‌کنشی:
مسیرها برنامه‌ریزی‌شده‌اند و باید از قطعی‌های کنترل‌گر جان سالم به در برند.
\end{solutionbox}

\begin{solutionbox}{\ref{exr:packetin}}
$10^5\times160$ بایت $= 16$ مگابایت بر ثانیه $= 128$ مگابیت بر ثانیه. با کل بسته‌ها: $10^5\times1532\times8 = 1.23$ گیگابیت بر
ثانیه -- ده برابر بیشتر، و کنترل‌گر باید آن‌ها را نیز ذخیره یا هدایت کند. بافرکردن در سوئیچ ظرفیت کانال کنترل را صرفه‌جویی
می‌کند.
\end{solutionbox}

\begin{solutionbox}{\ref{exr:little}}
$15\,000 = 1000\,(8 + T_{\mathrm{idle}}) \Rightarrow T_{\mathrm{idle}} = 7$ ثانیه. مهلت بیکاری بسیار کوتاه مدخل‌های جریان‌هایی را که
اندکی مکث می‌کنند بیرون می‌اندازد؛ بستهٔ بعدی آن‌ها دوباره تطابق نمی‌یابد و \lr{Packet-In} دیگری پدید می‌آورد که بار و تأخیر
کنترل‌گر را افزایش می‌دهد.
\end{solutionbox}

\begin{solutionbox}{\ref{exr:place}}
یک کنترل‌گر: فاصله‌های ۰، ۱، ۱، ۲، ۲، ۳؛ $\bar L = 9/6 = 1.5$ و $L_{\max} = 3$ میلی‌ثانیه. دو کنترل‌گر روبه‌رو: ۰، ۱، ۱، ۰، ۱، ۱؛
$\bar L = 4/6 = 0.67$ و $L_{\max} = 1$ میلی‌ثانیه. اگر یکی خراب شود، بازمانده با $L_{\max} = 3$ میلی‌ثانیه به همه خدمت می‌دهد --
جایابی باید برای حالت‌های خرابی هم ارزیابی شود.
\end{solutionbox}

\begin{solutionbox}{\ref{exr:consistent}}
اگر \lr{A} پیش از آنکه \lr{C} قاعده‌ای برای جریان داشته باشد به \lr{C} هدایت کند، \lr{C} تطابق نمی‌یابد و بسته را دور می‌ریزد (یا
به کنترل‌گر می‌فرستد). ترتیب مهم است: نخست روی \lr{C} نصب کن، سپس \lr{A} را تغییر بده («بساز پیش از آنکه بشکنی»). برای
به‌روزرسانی‌های کلی: قواعد جدید را روی \lr{C} و \lr{D} با تطابق برچسب نسخهٔ ۲ نصب کن (در حالی که قواعد نسخهٔ ۱ باقی‌اند)، سپس
\lr{A} را وادار کن بسته‌های جدید را با برچسب ۲ مهر کند؛ پس از تخلیهٔ بسته‌های قدیمی، قواعد نسخهٔ ۱ را حذف کن. به‌طور موقت
سوئیچ‌های متأثر هر دو نسخه را نگه می‌دارند -- تا دو برابر قواعد.
\end{solutionbox}

% ---------------------------------------------------------------------
\section{مطالعهٔ بیشتر}
مک‌کیون و همکاران \cite{mckeown2008openflow} \lr{OpenFlow} را معرفی کردند که مشخصات کنونی سوئیچ آن را \acro{ONF} منتشر
می‌کند \cite{onf2015openflow}؛ کروتز و همکاران \cite{kreutz2015sdn} \acro{SDN} را جامع مرور می‌کنند؛ \lr{RFC 7426}
\cite{rfc7426} لایه‌ها و اصطلاحات \acro{SDN} را تعریف می‌کند؛ باسهارت و همکاران \cite{bosshart2014p4} \lr{P4} را ارائه
می‌دهند. برای زمینهٔ موبایل \lr{TS 23.214} \cite{ts23214} و \lr{TS 29.244} \cite{ts29244} را ببینید.

\begin{latin}
\printbibliography[heading=subbibliography,title={References}]
\end{latin}
\end{document}
